% Generated draft: requires mathematical and source review.
\documentclass[11pt]{article}
\usepackage[margin=25mm]{geometry}
\usepackage{fontspec}
\setmainfont{Latin Modern Roman}
\usepackage{amsmath,amssymb,mathtools,unicode-math}
\setmathfont{Latin Modern Math}
\usepackage{enumitem,needspace,xurl,hyperref}
\setlength{\parindent}{0pt}
\setlength{\parskip}{5pt}
\setlength{\emergencystretch}{3em}
\newcommand{\N}{\mathbb{N}}
\newcommand{\Z}{\mathbb{Z}}
\newcommand{\Q}{\mathbb{Q}}
\newcommand{\R}{\mathbb{R}}
\newcommand{\C}{\mathbb{C}}
\newcommand{\F}{\mathbb{F}}
\newcommand{\A}{\mathbb{A}}
\newcommand{\PP}{\mathbb{P}}
\newcommand{\E}{\mathbb{E}}
\newcommand{\eps}{\varepsilon}
\newcommand{\dd}{\,\mathrm{d}}
\begin{document}
\section*{UnsolvedMath ID: 30002901}
\textbf{Source problem code:} catalogue entry
\subsection{Definitions and mathematical statement}
The supplied record does not specify the topology of the group, the completion denoted by \(C^*(G)\), or the assembly map. The following expresses its target under the discrete-group interpretation, while retaining the completion and assembly map as unresolved parameters.

Let \(G\) be a discrete group and let \(A=C^*(G)\) denote the intended group \(C^*\)-algebra. For a nonidentity element \(g\in G\) of finite order \(d\geq 2\), define
\[
p_g=\frac{1}{d}\sum_{k=1}^{d}g^k\in\Q G.
\]
Its image in \(A\) is a projection, and \([p_g]\in K_0(A)\) denotes its projection class. For this formulation, interpret the finite part as
\[
K_0^{\mathrm{fin}}(A)=\big\langle [p_g]:g\in G\setminus\{e\},\ \operatorname{ord}(g)<\infty\big\rangle_{\Z}\subseteq K_0(A).
\]
This generating convention is not explicitly defined in the supplied record. Let \(\mu:D\longrightarrow K_0(A)\) denote the intended assembly homomorphism; its domain \(D\) and construction remain unspecified.

The proposed assertion is that, for every group \(G\) in the intended scope, every integer \(n\geq 1\), and every choice of nonidentity finite-order elements \(g_1,\ldots,g_n\in G\) whose orders are pairwise distinct, the homomorphism
\[
\Phi:\Z^n\longrightarrow K_0^{\mathrm{fin}}(A),\qquad
(m_1,\ldots,m_n)\longmapsto\sum_{i=1}^{n}m_i[p_{g_i}]
\]
is injective and its image \(H\) satisfies
\[
H\cap\operatorname{im}(\mu)=\{0\}.
\]
Equivalently, these classes generate a free abelian subgroup of rank \(n\), and every nonzero element of that subgroup lies outside the assembly-map image. This is an unresolved statement-only formulation: identifying the completion, assembly map, and group hypotheses is necessary to determine a unique conjecture.
\subsection{Short English statement}
The question asks whether projections obtained by averaging over cyclic subgroups of distinct finite orders always give independent integral \(K\)-theory classes, with every nonzero integral combination outside the intended assembly-map image. The supplied record labels the problem open, but the group \(C^*\)-completion and assembly map require clarification.
\subsection{Sources}
\begin{itemize}
\item Supplied local problem registry record, permanent UnsolvedMath ID \(30002901\), exact source problem code \texttt{catalogue entry}, title ``Finite-Order Projections outside the Assembly-Map Image.'' The record gives proposed year 2015 and names no proposer.
\item Supplied report citation: \emph{Noncommutative Geometry (2016)}, Oberwolfach Reports, pp.~1631--1700, DOI \texttt{10.4171/owr/2015/29}. \url{https://doi.org/10.4171/owr/2015/29}. The supplied title year and DOI year differ; both are retained as recorded.
\item Supplied literature source: Weinberger--Yu, \emph{Geometry and Topology}. \url{https://msp.org/gt/2015/19-5/gt-v19-n5-p08-p.pdf}.
\item Supplied literature source: \emph{Higher rho invariants and the finite-part conjecture}. \url{https://doi.org/10.1016/j.aim.2016.11.030}.
\item Supplied literature source: ``Oberwolfach report 29.'' \url{https://doi.org/10.4171/owr/2015/29}.
\item Literature assessment dated \texttt{2026-08-21}: the supplied curation labels the problem open and reports that the finite-part \(K\)-theory conjecture is proved for broad classes such as finitely embeddable groups, while its general validity for arbitrary groups is not established by the report and subsequent literature. This is supplied dated curation, not an independently verified literature assessment.
\item The registry marks the statement \texttt{corrected\_verified} and reports a check against the official Oberwolfach report. That verification claim is supplied metadata; no source consultation or independent mathematical verification is asserted here.
\end{itemize}
\end{document}
